\section{Experimental Test Flow}\label{sec:test-flow}

The evaluation contrasts two mechanisms for autoscaling a sharded, stateful Redis
Cluster on the same six-node \texttt{k3s} cluster: (i) the native Kubernetes
\emph{Horizontal Pod Autoscaler} (HPA) acting on a self-managed cluster, and (ii) a
domain-specific \emph{operator} driven by \emph{Kubernetes Event-driven Autoscaling}
(KEDA). Both scenarios are exercised with an identical, deterministic six-phase
procedure so that the scaling mechanism is the sole independent variable:
(1)~scorch all prior state and scale the FastAPI entrypoint to zero;
(2)~provision the namespace, the Redis Cluster, and the entrypoint;
(3)~deploy the autoscaler;
(4)~drive a controlled \texttt{k6} workload and observe the scale-out;
(5)~quiesce the workload and perform a controlled scale-in; and
(6)~capture the resulting topology and key-space metrics in Grafana.
The workload, resource requests, monitoring stack, and autoscaler \texttt{behavior}
parameters are held constant; only the resharding intelligence differs. Bash scripts help to automate most ofthe procedure. For more details on the Bash scripts, please see the accompanying repository.

A Redis Cluster partitions its key space into $16{,}384$ hash slots distributed
across master nodes, so adding or removing a master requires \emph{resharding}---the
live migration of slots between nodes---rather than the mere addition of a pod
(\cite{redis:clusterspec,redis:scaling}). The central question of the test flow is
whether each mechanism can perform this slot-level work, because Kubernetes' native
workload primitives manage pod \emph{count} but possess no knowledge of cluster
\emph{topology}~(\cite{k8s:operator}).

\subsection{Scenario A: Native HPA on a Self-Managed Redis Cluster}\label{subsec:flow-hpa}

\begin{enumerate}
  \item \textbf{Provisioning.} A StatefulSet of six Redis pods is applied and the
  cluster is formed imperatively with
  \texttt{redis-cli -{}-cluster create -{}-cluster-replicas 1}, yielding three masters
  (\texttt{redis-0..2}) each with one replica (\texttt{redis-3..5}). Because a
  StatefulSet manifest cannot express the master--replica relationship, this role
  assignment is necessarily a runtime operation external to the declarative
  configuration~(\cite{k8s:statefulset,redis:clusterspec}).

  \item \textbf{Autoscaler.} An HPA targets the StatefulSet with a $40\%$ average
  CPU-utilisation goal; per-pod CPU requests are declared so the controller can
  compute utilisation~(\cite{k8s:hpa}). \texttt{minReplicas} is pinned to the six-pod
  baseline so the idle controller cannot scale down and strip the replicas before
  the run.

  \item \textbf{Scale-out.} The workload drives master CPU past the target. The HPA
  raises the desired replica count according to
  $\mathrm{desired} = \lceil \mathrm{current} \times
  (\mathrm{currentMetric}/\mathrm{targetMetric}) \rceil$~\cite{k8s:hpa} and appends
  pods beginning at \texttt{redis-6}. These pods join the StatefulSet but are never
  issued a \texttt{CLUSTER MEET} nor assigned slots, so each remains an isolated
  single-node island---a \emph{ghost pod} owning zero hash slots and serving no
  traffic.

  \item \textbf{Scale-in.} The HPA is removed and the StatefulSet is scaled down one
  pod at a time. StatefulSets terminate pods in reverse-ordinal order, from
  $\{N{-}1 \ldots 0\}$~(\cite{k8s:statefulset}), so the three replicas
  (\texttt{redis-5,4,3}) are deleted first; native cluster failover preserves every
  slot and no data is lost. The subsequent deletion removes a master
  (\texttt{redis-2}) whose replica is already gone. With no drain, no reshard, and no
  surviving replica to promote, that master's slots are orphaned and cluster-wide
  coverage falls below $16{,}384$.
\end{enumerate}

\paragraph{Hypothesis.} The HPA will provision capacity it structurally cannot
integrate on scale-out: ghost pods accumulate while throughput and latency fail to
recover, because the three original masters stay saturated. On scale-in it will
inflict irreversible data loss once redundancy is exhausted---the reverse-ordinal
\emph{data cliff}, observable as \texttt{cluster\_state:fail} and a step decrease in
total keys. Crucially, the replicas will \emph{delay} but not \emph{prevent} the
cliff, demonstrating that failover (protection against unplanned node death) is no
substitute for the slot migration a planned scale-in requires.

\subsection{Scenario B: Operator and KEDA on an Operator-Managed Redis Cluster}\label{subsec:flow-operator}

\begin{enumerate}
  \item \textbf{Provisioning.} The Opstree Redis Operator---a custom controller and
  \texttt{RedisCluster} custom resource that encode Redis-specific domain knowledge
  into a Kubernetes reconciliation loop~(\cite{k8s:operator,opstree:redisoperator})---is
  given a desired \texttt{clusterSize}. It declaratively creates the leader and
  follower StatefulSets and forms the cluster itself, issuing \texttt{CLUSTER MEET}
  and assigning slots, with no manual \texttt{redis-cli} step.

  \item \textbf{Autoscaler.} A KEDA \texttt{ScaledObject} watches a Prometheus
  cache-miss-ratio metric and actuates the operator's \texttt{clusterSize} through a
  managed HPA bound to the custom resource's \texttt{/scale}
  subresource~(\cite{keda:home}).

  \item \textbf{Scale-out.} When the metric crosses its threshold, KEDA raises
  \texttt{clusterSize}; the operator adds a master \emph{and} performs a
  topology-aware reshard, migrating a balanced slot range onto the new node so that
  the added capacity actually relieves load~(\cite{redis:scaling}).

  \item \textbf{Scale-in.} For reproducibility the cluster is first reset to a
  settled three-master baseline, then transitioned a single master at a time
  ($3\rightarrow4\rightarrow3$) via KEDA's \texttt{paused-replicas} annotation. On
  removal the operator drains the departing master---migrating its slots to the
  survivors---before deleting the pod. A documented limitation is captured: the
  operator has no self-healing for an interrupted reshard, which leaves a slot in the
  \texttt{importing}/\texttt{migrating} (``open'') state~(\cite{redis:clusterspec}) that
  must be repaired manually with \texttt{redis-cli -{}-cluster fix}.
\end{enumerate}

\paragraph{Hypothesis.} The operator will integrate each new master on scale-out
(full slot ownership and recovering throughput) and will drain gracefully on
scale-in, holding continuous $16{,}384$-slot coverage with zero data loss. This
safety is expected to come at the cost of substantially more provisioning machinery
than the HPA path and a residual need for manual intervention when a reshard is
interrupted---supporting the thesis that the operator is \emph{necessary but not, on
its own, sufficient} for fully hands-off elastic operation.

\subsection{The Necessity of Quiescing the Load Generator Before Scale-In}\label{subsec:quiescing}

Resharding is performed live and depends on clients following \texttt{ASK} and
\texttt{MOVED} redirections while a slot is held in the
\texttt{importing}/\texttt{migrating} state~(\cite{redis:clusterspec,redis:scaling}).
Under the write-heavy workload, a continuous stream of writes addressed to a slot
that is mid-migration is the dominant cause of an interrupted migration, leaving an
``open'' slot that the operator cannot repair on its own. Performing the controlled
scale-in while the generator is still firing therefore introduces a confound: an
interruption attributable to write pressure would be indistinguishable from a
deficiency of the scaling mechanism itself.

Because this study is scoped to the ghost-pod phenomenon and the operator's
resharding \emph{capability}---not its robustness under adversarial write load---the
controlled scale-in is conducted under quiescent conditions. The \texttt{k6}
generator is terminated and the procedure continues only once the request rate, and
the corresponding Redis command rate, have fallen to zero. This is a deliberate
experimental control rather than a concession for three reasons. First, it isolates
resharding correctness from a known confounding stressor. Second, it remains fair to
both scenarios, as the HPA data cliff is structural and manifests even at zero load,
so quiescing cannot conceal it. Third, it preserves measurability: the workload
configures no key expiration or eviction, so populated keys persist for post-hoc
accounting of any loss. The load-induced interruption is not discarded but reported
separately as a finding of the scale-out experiments.

\paragraph{Hypothesis.} Quiescing the generator before the scale-in will sharply
reduce the probability of an interrupted reshard, yielding a reproducible graceful
drain for the operator while leaving the HPA cliff unchanged. This isolates write
load---rather than the scaling mechanism alone---as the governing factor in reshard
interruption, and so permits a balanced comparison of operational complexity between
the two approaches.